% ============================================================================
% PODZIAŁ ZADAŃ I WKŁAD AUTORÓW
% ============================================================================
% Wymień konkretne rezultaty, nie tylko ogólne czynności. Udziały wszystkich
% autorów muszą sumować się do 100%. W projekcie indywidualnym wpisz jednego
% autora z udziałem 100%.
% ============================================================================

\chapter{Podział zadań i wkład autorów}
\label{chap:wklad-autorow}

Podział odpowiedzialności, wykonanych zadań i rezultatów przedstawiono
w \cref{tab:wklad-autorow}. Udział procentowy odnosi się do całego projektu,
a nie wyłącznie do objętości przygotowanego tekstu.

\begin{table}[htbp]
  \centering
  \caption{Podział zadań i udział autorów w realizacji projektu}
  \label{tab:wklad-autorow}
  \small
  \begin{tabularx}{\textwidth}{@{}p{2.8cm}X X p{1.5cm}@{}}
    \toprule
    Autor & Zadania i odpowiedzialność & Wytworzone rezultaty & Udział \\
    \midrule
    Imię i nazwisko
      & Np. opracowanie koncepcji, implementacja wybranego modułu i testy
      & Np. diagram architektury, kod modułu i przypadki testowe
      & \qty{100}{\percent} \\
    % Dla kolejnych autorów dodaj wiersze, a udziały zmień tak, aby suma
    % wynosiła 100%, np. 45% i 55%:
    % Imię i nazwisko & Zadania & Rezultaty & \qty{45}{\percent} \\
    \bottomrule
  \end{tabularx}
\end{table}

Autorzy potwierdzają zgodność powyższego zestawienia z rzeczywistym podziałem
pracy oraz akceptują przypisane zakresy odpowiedzialności i udziały procentowe.

% Historia repozytorium może potwierdzać przebieg pracy, ale sama liczba
% commitów nie jest miarą wartości wkładu poszczególnych osób.
